\chapter*{Введение}
\addcontentsline{toc}{chapter}{Введение}
\label{ch:Введение}

Обработка естественного языка на стороне клиента открывает возможность реализации интеллектуальных веб-приложений без передачи пользовательских данных на удалённый сервер. Мультиязычные трансформерные модели, конвертированные в формат \textit{ONNX} и исполняемые средствами \textit{ONNX Runtime Web} непосредственно в браузере, обеспечивают приватность, офлайн-функциональность и детерминированную задержку --- свойства, недостижимые при классическом клиент-серверном подходе. Вместе с тем однопоточная модель \textit{JavaScript}, ограничения памяти \textit{WebAssembly} и аппаратная неоднородность конечных устройств формируют комплекс ограничений, препятствующих прямому развёртыванию полноразмерных архитектур. Исследование методов квантования и дистилляции применительно к задаче анализа тональности на морфологически разнородных языках сохраняет высокую практическую значимость.

Актуальность работы: необходимость достижения приемлемого баланса между точностью классификации тональности и вычислительной эффективностью мультиязычных моделей глубокого обучения при их исполнении в браузерной среде. Устойчивость квантованных трансформерных архитектур к морфологической вариативности языков и их адаптация к ограничениям \textit{WebAssembly} и \textit{WebGPU} составляют центральный предмет настоящего исследования.

Цель выпускной квалификационной работы (ВКР): экспериментально оценить баланс между точностью и вычислительной эффективностью мультиязычных моделей глубокого обучения для задачи анализа тональности текста при их исполнении в браузерной среде с учётом квантования и аппаратных ограничений клиентских устройств.

Объект исследования: веб-технологии клиентского вывода нейросетевых моделей --- \textit{ONNX Runtime Web}, \textit{WebAssembly} и \textit{WebGPU} --- как среда исполнения мультиязычных трансформерных архитектур для задач обработки естественного языка.

Предмет исследования: точность классификации тональности и задержка вывода мультиязычных трансформерных моделей в зависимости от степени квантования (\textit{FP32} и \textit{INT8}) и морфологической структуры входного языка.

Задачи исследования:

\begin{enumerate}
    \item провести обзор современных методов обработки естественного языка и архитектур глубокого обучения с акцентом на их применимость в среде веб-браузера;
    \item исследовать методы оптимизации нейронных сетей --- дистилляцию знаний и квантование --- для адаптации ресурсоёмких моделей к ограничениям клиентских устройств;
    \item разработать архитектуру веб-приложения на базе \textit{React} с многопоточной обработкой данных через \textit{Web Workers};
    \item осуществить программную интеграцию \textit{ONNX Runtime Web} с поддержкой аппаратного ускорения \textit{WebGPU};
    \item выполнить дообучение выбранной архитектуры на трёхъязычном корпусе и провести сравнительное тестирование восьми конфигураций моделей;
    \item провести экспериментальную валидацию системы, оценив влияние квантования и морфологической структуры языка на задержку инференса и точность классификации.
\end{enumerate}

Теоретические основы исследования: в ходе анализа источников установлено, что тематика браузерного нейросетевого вывода и клиентской обработки естественного языка активно разрабатывается в зарубежной литературе. В работах Q.~Wang и соавторов, F.~Jia и соавторов детально исследованы характеристики бэкендов \textit{WebAssembly} и \textit{WebGPU} применительно к выводу моделей глубокого обучения. Методы квантования трансформерных архитектур рассмотрены в работах Z.~Yao и соавторов, A.~Bhandare и соавторов, а проблематика кросс-языкового переноса и мультиязычного анализа тональности --- в работах M.~Krasitskii и соавторов, T.~Filip и соавторов. Отечественные исследования (П.~В.~Стрюков и соавторы, З.~М.~Акиева и соавторы) охватывают методы квантования больших языковых моделей и применение \textit{WebAssembly} для ускорения вычислений в браузере. Вместе с тем сравнительная оценка влияния квантования на точность мультиязычного анализа тональности при браузерном исполнении остаётся недостаточно изученной.

Информационная база исследования: \textit{HuggingFace Hub}, \textit{Google Scholar}, \textit{eLIBRARY}, \textit{Amazon Reviews 2023}, \textit{RuReviews}, \textit{Kaggle}.

Методы исследования: дистилляция знаний, квантование, дообучение, компьютерный эксперимент, сравнительный анализ, бенчмаркирование.

Теоретическая значимость: систематизированы подходы к оптимизации трансформерных архитектур для браузерной среды исполнения; теоретически обоснованы критерии сравнения конфигураций квантования применительно к задачам мультиязычного анализа тональности; обобщены тенденции в проектировании конвейера клиентского \textit{NLP}-вывода с учётом морфологической структуры языка.

Практическая значимость: разработан и верифицирован полный конвейер клиентского вывода мультиязычных моделей тональности --- от дообучения до браузерного развёртывания, --- а полученные экспериментальные данные позволяют обоснованно выбирать архитектуру и степень квантования для конкретных сценариев применения.

Результаты исследования:

\begin{itemize}
    \item обоснована эффективность применения 8-битного квантования (\textit{INT8}) и \textit{SIMD}-инструкций для ускорения матричных операций в браузере, что позволяет достичь приемлемых показателей задержки при инференсе;
    \item выявлена нелинейная зависимость времени обработки от морфологического строя языка: агглютинативные и флективные языки (арабский, русский) создают большую нагрузку на токенизатор и память по сравнению с аналитическими (английский);
    \item подтверждена эффективность изоляции вычислительного процесса в среде \textit{Web Workers} для устранения блокировок интерфейса при выполнении ресурсоёмких тензорных операций;
    \item доказана практическая возможность полного отказа от \textit{Python}-бэкенда для задач классификации текста малого объёма посредством \textit{ONNX Runtime Web}.
\end{itemize}